\section{2010年高考题}
\subsection{2010年大纲I卷（理）}\label{gk:2010年大纲I卷（理）}

\subsection{2010年大纲I卷（文）}\label{gk:2010年大纲I卷（文）}

\subsection{2010年大纲II卷（理）}\label{gk:2010年大纲II卷（理）}

\subsection{2010年大纲II卷（文）}\label{gk:2010年大纲II卷（文）}

\subsection{2010年全国卷（理）}\label{gk:2010年全国卷（理）}

\subsection{2010年全国卷（文）}\label{gk:2010年全国卷（文）}

\subsection{2010年安徽卷（理）}\label{gk:2010年安徽卷（理）}

\subsection{2010年安徽卷（文）}\label{gk:2010年安徽卷（文）}

\subsection{2010年北京卷（理）}\label{gk:2010年北京卷（理）}

\begin{description}
    \item[一、]
    \begin{enumerate}
        \item 1
        \item 2
        \item 3
        \item 4
        \item 5
        \item 6
        \item 7
        \item 8
    \end{enumerate}
    \item[二、]
    \begin{enumerate}
      \addtocounter{enumi}{+8}   
        \item 1
        \item 2
        \item 3
        \item 4
        \item 1
        \item 6
    \end{enumerate}
    \item[三、]   
    \begin{enumerate}
      \addtocounter{enumi}{+14}   
        \item 1
        \item 2
        \item 3
        \item 已知函数\(f(x)=\ln(1+x)-x+\frac{k}{2}x^2 (k^2\geq 0)\).
        \begin{enumerate}[label=(\arabic*)]
            \item 当\(k=2\)时, 求曲线\(y=f(x)\)在店\((1,f(1))\)处的切线方程;
            \item 求\(f(x)\)的单调区间.
        \end{enumerate}
        \item 在平面直角坐标系 $xOy$ 中, 点 $B$ 与点 $A(-1, 1)$ 关于原点 $O$ 对称, $P$ 是动点, 且直线 $AP$ 与 $BP$ 的斜率之积等于 $-\frac{1}{3}$. 
    \begin{enumerate}[label=(\arabic*)]
        \item 求动点 $P$ 的轨迹方程;
        \item 设直线 $AP$ 和 $BP$ 分别与直线 $x = 3$ 交于点 $M$, $N$, 问: 是否存在点 $P$ 使得 $\triangle PAB$ 与 $\triangle PMN$ 的面积相等? 若存在, 求出点 $P$ 的坐标; 若不存在, 说明理由. 
    \end{enumerate}
        \item 已知集合 $S_{n} = \left\{ X \mid X = \left( x_{1}, x_{2}, \ldots, x_{n} \right), x_{i} \in \left\{ 0, 1 \right\}, i = 1, 2, \cdots, n \right\}$（$n \geqslant 2$）, 对于 $A = \left( a_{1}, a_{2}, \cdots, a_{n} \right), B = \left( b_{1}, b_{2}, \cdots, b_{n} \right) \in S_{n}$, 定义 $A$ 与 $B$ 的差为 $A - B = \left( \left| a_{1} - b_{1} \right|, \left| a_{2} - b_{2} \right|, \cdots, \left| a_{n} - b_{n} \right| \right)$;$A$ 与 $B$ 之间的距离为 $d \left( A, B \right) = \sum_{i=1}^{n} \left| a_{i} - b_{i} \right|$. 
\begin{enumerate}[label=(\arabic*)]
    \item 证明：$\forall A, B, C \in S_{n}$, 有 $A - B \in S_{n}$, 且 $d \left( A - C, B - C \right) = d \left( A, B \right)$;
    \item 证明：$\forall A, B, C \in S_{n}$, $d \left( A, B \right), d \left( A, C \right), d \left( B, C \right)$ 三个数中至少有一个是偶数;
    \item 设 $P \subseteq S_{n}$, $P$ 中有 $m$ $(m \geqslant 2)$ 个元素, 
    记 $P$ 中所有两元素间距离的平均值为 $\bar{d} \left( P \right)$. 
    证明：$\bar{d} \left( P \right) \leqslant \frac{mn}{2 \left( m - 1 \right)}$. 
\end{enumerate}
    \end{enumerate}
\end{description}
\subsection{2010年北京卷（文）}\label{gk:2010年北京卷（文）}

\begin{description}
    \item[一、]
    \begin{enumerate}
        \item 1
        \item 2
        \item 3
        \item 4
        \item 5
        \item 6
        \item 7
        \item 8
    \end{enumerate}
    \item[二、]
    \begin{enumerate}
      \addtocounter{enumi}{+8}   
        \item 1
        \item 2
        \item 3
        \item 4
        \item 1
        \item 6
    \end{enumerate}
    \item[三、]   
    \begin{enumerate}
      \addtocounter{enumi}{+14}   
        \item 1
        \item 2
        \item 3
        \item  设函数 $f(x)=\frac{a}{3}x^{3}+bx^{2}+cx+d(a>0)$, 且方程 $f'(x)-9x=0$ 的两个根分别为1, 4.
\begin{enumerate}[label=(\arabic*)]
    \item 当 $a=3$ 且曲线 $y=f(x)$ 过原点时, 求 $f(x)$ 的解析式; 
    \item 若 $f(x)$ 在 $(-\infty,+\infty)$ 内无极值点, 求 $a$ 的取值范围.
\end{enumerate}
        \item 已知椭圆C的左、右焦点坐标分别是 $(-\sqrt{2},0)$, $(\sqrt{2},0)$, 离心率是 $\frac{\sqrt{6}}{3}$.直线 $y=t$ 与椭圆C交于不同的两点 $M$, $N$, 以线段 $MN$ 为直径作圆 $P$, 圆心为 $P$.
\begin{enumerate}[label=(\arabic*)]
    \item 求椭圆 $C$ 的方程; 
    \item 若圆 $P$ 与 $x$ 轴相切, 求圆心 $P$ 的坐标; 
    \item 设 $Q(x,y)$ 是圆 $P$ 上的动点, 当 $t$ 变化时, 求 $y$ 的最大值.
\end{enumerate}
        \item 已知集合 $S_{n} = \left\{ X \mid X = \left( x_{1}, x_{2}, \ldots, x_{n} \right), x_{i} \in \left\{ 0, 1 \right\}, i = 1, 2, \cdots, n \right\}$（$n \geqslant 2$）, 
        对于 $A = \left( a_{1}, a_{2}, \cdots, a_{n} \right), B = \left( b_{1}, b_{2}, \cdots, b_{n} \right) \in S_{n}$,
         定义 $A$ 与 $B$ 的差为 $A - B = \left( \left| a_{1} - b_{1} \right|, \left| a_{2} - b_{2} \right|,\right.$ $ \cdots, $ $\left. \left| a_{n} - b_{n} \right| \right)$; $A$ 与 $B$ 之间的距离为 $d \left( A, B \right) = \sum_{i=1}^{n} \left| a_{i} - b_{i} \right|$.
\begin{enumerate}[label=(\arabic*)]
    \item 当 $n=5$ 时, 设 $A=\left(0,1,0,0,1\right)$, $B=\left(1,1,1,0,0\right)$, 求 $A-B$ 和 $d(A,B)$; 
    \item 证明：$\forall A, B, C \in S_{n}$, 有 $A - B \in S_{n}$, 且 $d \left( A - C, B - C \right) = d \left( A, B \right)$; 
    \item 证明：$\forall A, B, C \in S_{n}$, $d \left( A, B \right), d \left( A, C \right), d \left( B, C \right)$ 三个数中至少有一个是偶数.
\end{enumerate}
    \end{enumerate}
\end{description}
\subsection{2010年重庆卷（理）}\label{gk:2010年重庆卷（理）}

\subsection{2010年重庆卷（文）}\label{gk:2010年重庆卷（文）}

\subsection{2010年福建卷（理）}\label{gk:2010年福建卷（理）}

\subsection{2010年福建卷（文）}\label{gk:2010年福建卷（文）}

\subsection{2010年广东卷（理）}\label{gk:2010年广东卷（理）}

\subsection{2010年广东卷（文）}\label{gk:2010年广东卷（文）}

\subsection{2010年湖北卷（理）}\label{gk:2010年湖北卷（理）}

\subsection{2010年湖北卷（文）}\label{gk:2010年湖北卷（文）}

\subsection{2010年湖南卷（理）}\label{gk:2010年湖南卷（理）}

\subsection{2010年湖南卷（文）}\label{gk:2010年湖南卷（文）}

\subsection{2010年江苏卷}\label{gk:2010年江苏卷}

\subsection{2010年江西卷（理）}\label{gk:2010年江西卷（理）}

\subsection{2010年江西卷（文）}\label{gk:2010年江西卷（文）}

\subsection{2010年辽宁卷（理）}\label{gk:2010年辽宁卷（理）}

\subsection{2010年辽宁卷（文）}\label{gk:2010年辽宁卷（文）}

\subsection{2010年上海卷（理）}\label{gk:2010年上海卷（理）}

\subsection{2010年上海卷（文）}\label{gk:2010年上海卷（文）}

\subsection{2010年四川卷（理）}\label{gk:2010年四川卷（理）}

\subsection{2010年四川卷（文）}\label{gk:2010年四川卷（文）}

\subsection{2010年山东卷（理）}\label{gk:2010年山东卷（理）}

\subsection{2010年山东卷（文）}\label{gk:2010年山东卷（文）}

\subsection{2010年陕西卷（理）}\label{gk:2010年陕西卷（理）}

\subsection{2010年陕西卷（文）}\label{gk:2010年陕西卷（文）}

\subsection{2010年天津卷（理）}\label{gk:2010年天津卷（理）}

\subsection{2010年天津卷（文）}\label{gk:2010年天津卷（文）}

\subsection{2010年浙江卷（理）}\label{gk:2010年浙江卷（理）}

\subsection{2010年浙江卷（文）}\label{gk:2010年浙江卷（文）}






